% !TeX root = ../main.tex
\section{A simple feedforward network}

In this task, we simulate a network of neurons in an aim to achieve complex behavior.
We connect several populations of neurons with each other, such that they form a chain.
After an initial excitation, each population in the chain spikes and then excites the next population in the chain.
These are the \quotedef{excitatory} populations.
For each excitatory population, we introduce an \quotedef{inhibitory} population which gets excited by the same (previous) population.
After spiking, it inhibits its accompanying excitatory population, thereby hindering it from spiking again.
In this way, the initial stimulation travels through this so-called \quotedef{synfire chain}.

The notebook corresponding to this task was pre-filled with an implementation of a synfire chain.
In the following, we use this implementation to answer several questions about the inner workings of a feedforward network.
Specifically useful for quick inspection are the two plots which the notebook produced for each run.
The first subplot is a plot of the pair (time of spike, neuron id) for each spiked neuron.
The second subplot is the membrane trace of the first excited neuron.

\subsection{Adjusting Weights}
We are asked to fill in the weights which determine how much one population excites or inhibits another.
There are five weights: \verb|stim_exc|, \verb|stim_inh|, \verb|exc_exc|, \verb|exc_inh|, and \verb|inh_exc|.
The allowed values are integers from 0-63 (negative in the last case).

